% Part 3: Proposed TA-PPO Framework

\section{Traffic-Aware Proximal Policy Optimization}

This section presents our Traffic-Aware PPO (TA-PPO) framework, which integrates traffic prediction with deep reinforcement learning for proactive routing.

\subsection{Architecture Overview}

Figure \ref{fig:architecture} illustrates the complete TA-PPO architecture, consisting of four main components:

\begin{enumerate}
\item \textbf{Data Collector}: Continuously monitors and buffers link metrics over a sliding window
\item \textbf{Prediction Module}: LSTM/GRU networks forecast future congestion
\item \textbf{PPO Agent}: Actor-critic networks learn optimal routing policy
\item \textbf{Routing Executor}: Implements routing decisions and updates network state
\end{enumerate}

% Include architecture diagram
\begin{figure*}[t]
\centering
\includegraphics[width=0.9\textwidth]{architecture_diagram.png}
\caption{Traffic-Aware PPO Architecture. The system integrates data collection, traffic prediction, and DRL-based routing in a closed loop. Historical metrics feed into LSTM/GRU models to predict future congestion, which enhances the state representation for the PPO agent.}
\label{fig:architecture}
\end{figure*}

\subsection{Traffic Prediction Module}

\subsubsection{Data Collection}

The Data Collector maintains a sliding window buffer for each ISL, tracking seven key metrics:

\begin{equation}
\mathbf{H}_{ij}(t) = [h_{ij}(t-W+1), ..., h_{ij}(t)]
\end{equation}

where $W$ is the lookback window (typically 50 timesteps) and each $h_{ij}(\tau)$ contains:

\begin{equation}
h_{ij}(\tau) = [Q_{ij}(\tau), U_{ij}(\tau), B_{ij}(\tau), D_{ij}(\tau), \delta_{ij}(\tau), \rho_{ij}(\tau), SR_{ij}(\tau)]
\end{equation}

representing queue length, utilization, bandwidth used, packet drops, average delay, packet arrival rate, and success rate respectively.

\subsubsection{LSTM Prediction Model}

The LSTM architecture consists of:

\begin{equation}
\begin{split}
\mathbf{f}_t &= \sigma(\mathbf{W}_f \cdot [\mathbf{h}_{t-1}, \mathbf{x}_t] + \mathbf{b}_f) \\
\mathbf{i}_t &= \sigma(\mathbf{W}_i \cdot [\mathbf{h}_{t-1}, \mathbf{x}_t] + \mathbf{b}_i) \\
\mathbf{o}_t &= \sigma(\mathbf{W}_o \cdot [\mathbf{h}_{t-1}, \mathbf{x}_t] + \mathbf{b}_o) \\
\tilde{\mathbf{C}}_t &= \tanh(\mathbf{W}_C \cdot [\mathbf{h}_{t-1}, \mathbf{x}_t] + \mathbf{b}_C) \\
\mathbf{C}_t &= \mathbf{f}_t \odot \mathbf{C}_{t-1} + \mathbf{i}_t \odot \tilde{\mathbf{C}}_t \\
\mathbf{h}_t &= \mathbf{o}_t \odot \tanh(\mathbf{C}_t)
\end{split}
\end{equation}

where $\mathbf{f}_t, \mathbf{i}_t, \mathbf{o}_t$ are forget, input, and output gates; $\mathbf{C}_t$ is the cell state; $\sigma$ is the sigmoid function; and $\odot$ denotes element-wise multiplication.

The prediction layer maps the final hidden state to future predictions:

\begin{equation}
\hat{Q}_{ij}(t+1:t+H) = \mathbf{W}_{\text{out}} \cdot \mathbf{h}_t + \mathbf{b}_{\text{out}}
\end{equation}

where $H$ is the prediction horizon (typically 10 timesteps).

\subsubsection{GRU Prediction Model}

The GRU architecture simplifies LSTM by combining forget and input gates:

\begin{equation}
\begin{split}
\mathbf{z}_t &= \sigma(\mathbf{W}_z \cdot [\mathbf{h}_{t-1}, \mathbf{x}_t]) \\
\mathbf{r}_t &= \sigma(\mathbf{W}_r \cdot [\mathbf{h}_{t-1}, \mathbf{x}_t]) \\
\tilde{\mathbf{h}}_t &= \tanh(\mathbf{W} \cdot [\mathbf{r}_t \odot \mathbf{h}_{t-1}, \mathbf{x}_t]) \\
\mathbf{h}_t &= (1 - \mathbf{z}_t) \odot \mathbf{h}_{t-1} + \mathbf{z}_t \odot \tilde{\mathbf{h}}_t
\end{split}
\end{equation}

where $\mathbf{z}_t$ is the update gate and $\mathbf{r}_t$ is the reset gate.

GRU has fewer parameters than LSTM, often resulting in faster training and inference while maintaining comparable accuracy.

\subsubsection{Training Procedure}

The prediction models are trained offline using collected traffic data. The loss function is Mean Squared Error (MSE):

\begin{equation}
\mathcal{L}_{\text{pred}} = \frac{1}{N \cdot H} \sum_{n=1}^{N} \sum_{h=1}^{H} \left(Q_{ij}(t+h) - \hat{Q}_{ij}(t+h)\right)^2
\end{equation}

We use Adam optimizer with learning rate $\eta = 0.001$ and early stopping with patience of 10 epochs. Data is normalized using:

\begin{equation}
\mathbf{x}_{\text{norm}} = \frac{\mathbf{x} - \mu}{\sigma}
\end{equation}

where $\mu$ and $\sigma$ are mean and standard deviation computed from training data.

\subsection{PPO Agent}

\subsubsection{Actor-Critic Architecture}

The PPO agent consists of two neural networks sharing a common feature extractor:

\textbf{Feature Extractor}:
\begin{equation}
\mathbf{z}_t = \text{ReLU}(\mathbf{W}_1 \cdot \mathbf{s}_t + \mathbf{b}_1)
\end{equation}

\textbf{Actor (Policy) Network}:
\begin{equation}
\pi_\theta(a_t | \mathbf{s}_t) = \text{Softmax}(\mathbf{W}_\pi \cdot \mathbf{z}_t + \mathbf{b}_\pi)
\end{equation}

\textbf{Critic (Value) Network}:
\begin{equation}
V_\phi(\mathbf{s}_t) = \mathbf{W}_V \cdot \mathbf{z}_t + \mathbf{b}_V
\end{equation}

\subsubsection{PPO Objective}

The PPO objective maximizes:

\begin{equation}
\mathcal{L}^{\text{PPO}}(\theta) = \mathbb{E}_t\left[\min\left(r_t(\theta)\hat{A}_t, \text{clip}(r_t(\theta), 1-\epsilon, 1+\epsilon)\hat{A}_t\right)\right]
\end{equation}

where:
\begin{itemize}
\item $r_t(\theta) = \frac{\pi_\theta(a_t|\mathbf{s}_t)}{\pi_{\theta_{\text{old}}}(a_t|\mathbf{s}_t)}$ is the probability ratio
\item $\hat{A}_t$ is the advantage estimate computed using Generalized Advantage Estimation (GAE):
\begin{equation}
\hat{A}_t = \sum_{l=0}^{\infty} (\gamma \lambda)^l \delta_{t+l}
\end{equation}
where $\delta_t = R_t + \gamma V(\mathbf{s}_{t+1}) - V(\mathbf{s}_t)$
\item $\epsilon$ is the clipping parameter (typically 0.15)
\end{itemize}

The complete loss function combines policy, value, and entropy terms:

\begin{equation}
\mathcal{L}(\theta) = \mathcal{L}^{\text{PPO}}(\theta) - c_1 \mathcal{L}^V(\theta) + c_2 \mathcal{H}(\pi_\theta)
\end{equation}

where:
\begin{itemize}
\item $\mathcal{L}^V(\theta) = (V_\theta(\mathbf{s}_t) - V_{\text{target}})^2$ is the value loss
\item $\mathcal{H}(\pi_\theta) = -\sum_a \pi_\theta(a|\mathbf{s}) \log \pi_\theta(a|\mathbf{s})$ is the entropy bonus
\item $c_1 = 0.7$ and $c_2 = 0.005$ are coefficients
\end{itemize}

\subsection{Predictive Routing Cost Function}

The key innovation of TA-PPO is incorporating predicted congestion into routing decisions. The routing cost for link $(i,j)$ is:

\begin{equation}
C_{ij}(t) = w_1 \tau_{ij}(t) + w_2 Q_{ij}(t) + w_3 \hat{Q}_{ij}(t+h)
\end{equation}

where:
\begin{itemize}
\item $w_1 = 1.0$ weights propagation delay
\item $w_2 = 2.0$ weights current congestion
\item $w_3 = 5.0$ weights predicted congestion (higher to prioritize avoidance)
\end{itemize}

This cost function enables proactive routing: even if a link is currently uncongested ($Q_{ij}(t) \approx 0$), it will be avoided if predicted to become congested ($\hat{Q}_{ij}(t+h) > 0$).

\subsection{Integration and Execution Flow}

Algorithm \ref{alg:tappo} presents the complete TA-PPO execution flow.

\begin{algorithm}[t]
\caption{Traffic-Aware PPO Routing}
\label{alg:tappo}
\begin{algorithmic}[1]
\STATE Initialize PPO networks $\pi_\theta$, $V_\phi$
\STATE Load pre-trained LSTM/GRU model $M_{\text{pred}}$
\STATE Initialize data collector with buffer size $W$
\FOR{each episode}
    \STATE Reset network state
    \FOR{each timestep $t$}
        \FOR{each packet at satellite $s_i$}
            \STATE Collect current metrics for all neighbors
            \IF{$t \geq W_{\text{warmup}}$}
                \STATE Retrieve history $\mathbf{H}_{ij}$ for each neighbor $j$
                \STATE Predict: $\hat{Q}_{ij}(t+h) = M_{\text{pred}}(\mathbf{H}_{ij})$
            \ELSE
                \STATE Set $\hat{Q}_{ij}(t+h) = 0$ (no prediction)
            \ENDIF
            \STATE Construct state $\mathbf{s}_t$ with predictions
            \STATE Sample action: $a_t \sim \pi_\theta(\cdot | \mathbf{s}_t)$
            \STATE Forward packet to selected neighbor
            \STATE Observe reward $R_t$ and next state $\mathbf{s}_{t+1}$
            \STATE Store $(\mathbf{s}_t, a_t, R_t, \mathbf{s}_{t+1})$ in buffer
        \ENDFOR
        \STATE Update network topology $\mathcal{G}(t+1)$
        \STATE Update data collector with new metrics
    \ENDFOR
    \IF{buffer contains $\geq B$ samples}
        \STATE Compute advantages using GAE
        \FOR{$K$ epochs}
            \STATE Sample mini-batch from buffer
            \STATE Update $\theta$ by maximizing $\mathcal{L}(\theta)$
        \ENDFOR
    \ENDIF
\ENDFOR
\end{algorithmic}
\end{algorithm}

The algorithm operates in three phases:

\begin{enumerate}
\item \textbf{Warmup Phase} ($t < W_{\text{warmup}}$): Collect historical data without predictions, using baseline routing.

\item \textbf{Prediction Phase} ($t \geq W_{\text{warmup}}$): Use trained prediction models to forecast congestion and enhance state representation.

\item \textbf{Learning Phase}: Continuously update PPO networks using collected experience.
\end{enumerate}

\subsection{Computational Complexity}

The computational complexity per routing decision is:

\begin{itemize}
\item Prediction: $O(W \cdot H \cdot d_h)$ for LSTM/GRU, where $d_h$ is hidden dimension
\item Policy evaluation: $O(N_{\max} \cdot d_f + d_f \cdot d_\pi)$ for actor network
\item Total: $O(W \cdot H \cdot d_h + N_{\max} \cdot d_f)$
\end{itemize}

With $W=50$, $H=10$, $d_h=128$, $N_{\max}=8$, $d_f=256$, this is tractable for onboard satellite processors.

